%---------------------------------------------------------------------
\chapter{Introduction}
\label{ch:intro}
%---------------------------------------------------------------------

\section{Background}
India generates large quantities of electrical and electronic waste, and the E-Waste (Management)
Rules, 2022 \cite{ewaste2022}, notified under the Environment (Protection) Act, 1986, regulate how that
waste is collected and recycled. The Rules place an Extended Producer Responsibility (EPR) on every
producer of electrical and electronic equipment (EEE). A producer does not recycle its own products. It
instead buys EPR certificates from registered recyclers, and each certificate represents recovered
material of a given metal. The Central Pollution Control Board (CPCB) administers the mechanism through an
online portal, on which registrations, obligations, certificate generation, certificate transfers and
filed returns are recorded.

For a compliance officer in a producer firm, the central question each period is whether the obligation
for each metal has been met by credits that are genuinely usable. The answer is never a single number
read from the portal. The obligation comes from the Rules and from the producer's declared sales. The
credits come from certificates, which can be issued, transferred, cancelled or expired. A payment to a
recycler is a commercial fact and is not itself a credit, because the credit exists only after the
portal transfer is completed. A portal balance is a snapshot taken at one time and can be stale. Each
of these facts lives in a different document or system, so the officer reconciles them by hand across
invoices, bank records, ledger exports, portal statements and the text of the Rules.

\section{Problem Statement}
Manual reconciliation of EPR obligations is slow and error-prone, and the errors are of a specific kind.
Overlapping exports count a transfer twice. Units are mixed between kilograms and tonnes. An invoice is
issued under an entity alias that is not the registered recycler. A certificate that was cancelled or
expired is still counted. A recycler registered for one metal issues credits for another. Credits in the
current year are used to clear arrears from an earlier year. A memo is read as an official extension of a
deadline although no amendment exists. An apparent shortfall can therefore come from a record error, from
a pending transfer, from an inapplicable rule, or from a real unmet obligation, and each calls for a
different action.

The decision that this project supports is stated precisely:

\begin{quote}
\emph{Which discrepancy should the compliance officer investigate next, and what verification or
corrective action does the evidence support?}
\end{quote}

The output is an advisory working paper for a human reviewer. It is not a regulatory filing, a legal
certificate of compliance, a finding of fraud, or an assessment of a penalty, and the system is built
so that it cannot produce those outputs.

\section{Motivation and Why Agentic AI}
A fixed script can recompute an obligation balance, and the project does exactly that with deterministic
code. A script cannot decide what to look at next. When a gap appears, the useful next step depends on
what is found. A gap with a pending transfer sends the officer to the transfer record. A gap with a
certificate marked cancelled sends the officer to the certificate history. A rule that seems to apply
sends the officer to the text of the Rules and to the question whether that version was in force for
the period. The number of useful steps varies from case to case, and a stopping rule is needed when
nothing new can be found.

A single prompt to a language model has the opposite weakness. It can produce fluent text, but it cannot be
trusted to calculate exactly, to cite a provision that really says what it claims, or to refuse to guess.
For a compliance setting the cost of a confident but unsupported statement is high.

An agent that plans a check, calls a tool, inspects the observed result and then revises or stops
fits this problem. The design in this report keeps the useful part of the agent, which is the choice
of the next check, and removes the parts that a language model should not own. Arithmetic, date handling
and citation verification are deterministic. The model proposes, and bounded code decides what is accepted.
The project therefore treats the model as an investigator with a narrow, audited interface and not as an
oracle.

\section{Objectives}
The objectives of the project are the following.
\begin{enumerate}[itemsep=2pt]
\item Build a local agentic system that, for one producer and one reporting period, plans reconciliation
      checks, calls audited tools, assesses each observation, and either continues, resolves or abstains
      within fixed budgets.
\item Keep every quantity exact: obligations, supported allocations and gaps are computed with decimal
      arithmetic by deterministic tools and are never produced by the model.
\item Ground regulatory statements in retrieved provisions with source, version and locator, and prefer
      an explicit abstention to an unsupported claim.
\item Record every plan, tool call, model call, observation and stop reason in an append-only audit trail
      from which a run can be inspected and replayed.
\item Require a human reviewer to decide each report, bound to the exact report and input versions.
\item Evaluate the system on a synthetic benchmark with an independent reference, comparing fixed
      retrieval with bounded adaptive investigation, and report all results, including negative ones.
\end{enumerate}

\section{Scope of the Project}
\label{sec:scope}
The scope is deliberately narrow so that the system can be finished and evaluated.
\begin{itemize}[itemsep=2pt]
\item \textbf{In scope:} one producer and one liability year per case; the metals copper, aluminium, iron
      and gold as modelled in the calculation profile; structured and semi-structured evidence (PDF,
      CSV, JSON); a curated regulatory corpus; synthetic dossiers; a local model served on the same machine;
      a command-line interface, an HTTP API and a browser interface for review.
\item \textbf{Out of scope:} live portal integration; regulatory filing; any legal determination of
      compliance; inference of intent or fraud; calculation of compensation or penalties; credit trading;
      real producer data; and cloud model providers. Solar photovoltaic waste and non-notified regimes are
      outside the calculation profile, and the system is expected to say so rather than to compute them.
\end{itemize}
All dossiers are fictional and are watermarked as synthetic. No result in this report should be read as a
statement about the behaviour of any real producer or recycler.

\section{Organisation of the Report}
Chapter~\ref{ch:lit} reviews existing practice and related technical work and states the gap.
Chapter~\ref{ch:req} derives the requirements and the data. Chapter~\ref{ch:design} presents the
architecture, components and agentic patterns. Chapter~\ref{ch:concepts} maps the project to the agentic
concepts required by the course. Chapter~\ref{ch:impl} describes the implementation.
Chapter~\ref{ch:eval} covers observability and evaluation, and Chapter~\ref{ch:exp} the controlled
experiment. Chapter~\ref{ch:limits} discusses limitations, ethics and safety, and Chapter~\ref{ch:concl}
concludes. The appendices give the contribution of each team member, run instructions, annotated traces
and the prompts.

\pending{Chapters~\ref{ch:eval} and~\ref{ch:exp} contain the held-out results. The held-out runs have not
been made, so those tables are empty and are marked in place. Every number elsewhere in this report that is
tagged as a development result comes from runs on the 24 development dossiers.}
